\section{Caching in Link Traversal-based Query Processing}
To validate the query sequences generated by the SolidBenchSequence (SBS) benchmark,
we implemented three caching approaches in the Comunica link traversal engine~\cite{taelman2018comunica,taelman2023link}.
Specifically, we cache the quads \rt{I would call them "triples" here, because formally, "quad" is not a formal concept in the RDF 1.1 spec. (it will be in 1.2). You could optionally add a footnote for clarity saying that you also capture named graph info.} from documents dereferenced during link traversal.
High cache hit rates eliminate redundant HTTP requests, increasing data availability \rt{The use of the word "availability" can be confusing here. Readers may confuse it with caches leading to higher server availability due to less requests; but I don't think this is what you're trying to say?} and reducing query latency. 
To evaluate how different cache management strategies impact this latency,
we compare the following three configurations:

\rt{At this point, I realized that you never actually explain how LTQP works (link queue, reachability semantics, ...). Some minimal explanation of this must be added, to make this work self-contained, to make it possible for readers not knowing about LTQP to understand this paper. For example add this in the introduction or literature review.}

\paragraph{Unindexed Cache}
This baseline approach caches the raw \rt{Unclear what "raw" refers to here. Do you mean "in-memory"?} quads obtained during traversal.
Upon a cache hit, these quads are retrieved and indexed into the active local store \rt{Unclear to the reader what the "active local store" is. This for example should be explained in the general description of LTQP.} for query execution.
\rt{Do I understand it correctly that this approach is essentially just an array of quads/triples?}

\paragraph{Indexed Cache}
Alternatively, we index quads before caching them. 
Upon a cache hit, the engine extracts only quads matching the quad patterns in the current query.
While this drastically reduces the memory footprint and indexing overhead of the local store \rt{It's unclear why this reduces memory usage. The same data is stored, right? Just in a different way?},
it increases the initial cache ingestion cost. 
We hypothesize this trade-off is only beneficial given sufficiently high hit rates.

\paragraph{Indexed Cache-based Estimation}
\rt{To clarify that this is really an extension of Indexed Cache, perhaps we could call it "Indexed Cache + Cardinality estimation"?}
Because the cache is indexed, it can execute count queries to estimate quad pattern cardinalities.
Traditional LTQP relies on heuristics due to a lack of prior knowledge~\cite{hartig2011zero}.
By estimating cardinalities against \emph{all} cached documents,
regardless of their reachability for the current query,
we can replace these heuristics with Comunica's default cost-based optimizer~\cite{taelman2018comunica} \rt{In link traversal default, we use Olaf's zero-knowledge heuristics, so it'd be good to distinguish to this approach as well.}
to generate more accurate query plans. \\

\rt{It may be helpful for the reader to visually represent these 3 approaches next to each other.}

For each caching approach, we investigate three sizes \rt{I would emphasize here again that you focus on in-memory storage}: small (350,000 quads), medium (700,000 quads),
and large (1,400,000 quads). These triple counts correspond to about 10\%, 20\%, and 40\% of the total
generated dataset size.

Our implementations adhere to the RFC-9111 specification~\cite{fielding2022rfc} \rt{Good! But what does that mean? (most readers won't know this RFC, and won't care to look it up)}.
However, because the benchmark dataset is static, cache entries require no revalidation and never become stale. 
Consequently, the observed hit rates are likely higher than in a live environment. 
Furthermore, while the fundamental concept of caching in link traversal is established~\cite{hartig2011caching} \rt{This hartig2011caching paper should really be discussed in more detail. Perhaps you can add it to this dedicated LTQP description I asked for above?},
integrating and comparing these distinct caching strategies provides novel and strong baselines
for evaluating caching in traversal engines, upon which future research
into advanced caching approach can build.